\scanpage{29}
\mspage{27}
% archive sheet number: 28
s'identifie à
\[
\mathbb D_{\underline{\mathrm{Hom}}(I^{\unsure{\circ|'}},\mathcal M)}=\mathbb D_{\mathcal M(I)}
\]
i.e.
\[
(\mathbb D_{\mathcal M})_I\unsure{=|\simeq}\mathbb D_{\mathcal M(I)}
\]

\deleted{\unsure{Et si|Soit donc}
\[
I\xrightarrow{f}\overset{\unsure{\longleftrightarrow}}{I'}
\]
les foncteurs $f^*_{\mathbb D}$, $f_*^{\mathbb D}$, $f_!^{\mathbb D}$ entre $\mathbb D_I$ et $\mathbb D_{I'}$ s'interprètent comme des foncteurs entre $\mathbb D_{\mathcal M(I)}$ et $\mathbb D_{\mathcal M(I')}$.}

\uline{NB} Ici $\mathcal M$ et $\mathbb D_{\mathcal M}$ se déterminent mutuellement -- le foncteur
\[
\begin{aligned}
\mathcal M&\unsure{\mathord{-}|\longrightarrow}\mathbb D_{\mathcal M}\\
(\mathrm{Cat})&\longrightarrow\mathrm{DER}
\end{aligned}
\]
est 2-fidèle, i.e. si $\mathcal M,\mathcal M'$ sont \unsure{dans|deux} $\mathrm{Cat}$
\[
\underline{\mathrm{Hom}}(\mathcal M,\mathcal M')\xrightarrow{\approx}\mathrm{Hom}(\mathbb D_{\mathcal M},\mathbb D_{\mathcal M'})
\]
est une équivalence de catégories. D'autre part, on vérifie tout de suite que \deleted{$\unsure{f_I}$} \deleted{$\unsure{f}$} \inserted{$u$}$:\mathcal M\longrightarrow\mathcal M'$ est exact à g (resp. à droite) ssi. $\mathbb D_u:\mathbb D_{\mathcal M}\longrightarrow\mathbb D_{\mathcal M'}$ l'est.

Ainsi, pour voir si un \deleted{foncteur} \inserted{morphisme} de dérivateurs de la forme
\[
f^{\mathbb D}_*:\deleted{\unsure{\mathcal M_*}}\,
\underset{\substack{\unsure{\shortparallel|\simeq}\\\mathbb D_{\mathcal M(I)}}}{(\mathbb D_{\mathcal M}\deleted{\unsure{f}})_I}
\longrightarrow
\underset{\substack{\shortparallel\\\mathbb D_{\mathcal M(I')}}}{(\mathbb D_{\mathcal M})_{I'}}
\]
est exact à \uline{droite}, il revient au même de dire
